% =====================================================================
% FONTES DEPENDENTES EM MALHAS --- NIVEL 2
% Questoes integradas ao bloco de Aplicacao do Capitulo 9.
% Correntes de malha desenhadas no sentido HORARIO.
% =====================================================================

\begin{questao}[2]{}
No circuito, a fonte de tensão dependente fornece $2{,}5I_1$ volts e atua na
malha 2. Determine $I_1$, $I_2$ e a corrente real no resistor central.

\circuitoq{%
\begin{circuitikz}[american,scale=.92,transform shape,line width=.8pt]
\draw (0,0) to[\fonteV,l^={$\SI{20}{\volt}$}] (0,3.2)
      to[R,l={$\SI{5}{\ohm}$}] (3.2,3.2) coordinate(t);
\draw (t) to[R,l={$\SI{5}{\ohm}$}] (3.2,0) coordinate(b) -- (0,0);
\draw (t) to[R,l={$\SI{5}{\ohm}$}] (6.4,3.2)
      to[cvsource,l={$2{,}5I_1$}] (6.4,0) -- (b);
\draw[->,loopcol,line width=1.2pt] (1.75,1.00) arc[start angle=0,end angle=-310,radius=.46cm];
\node[loopcol,fill=white,inner sep=1pt] at (1.75,1.00){\footnotesize$I_1$};
\draw[->,loopcol,line width=1.2pt] (4.75,1.00) arc[start angle=0,end angle=-310,radius=.46cm];
\node[loopcol,fill=white,inner sep=1pt] at (4.75,1.00){\footnotesize$I_2$};
\end{circuitikz}}

\resposta{$I_1=\SI{3.2}{\ampere}$; $I_2=\SI{2.4}{\ampere}$; $I_c=\SI{0.8}{\ampere}$}
\resolucao{Com correntes de malha horárias,
\[
10I_1-5I_2=20,
\qquad
-5I_1+10I_2=2{,}5I_1.
\]
A segunda equação dá $I_2=0{,}75I_1$. Substituindo na primeira,
$I_1=\SI{3.2}{\ampere}$ e $I_2=\SI{2.4}{\ampere}$. No ramo comum,
$I_c=I_1-I_2=\SI{0.8}{\ampere}$.}
\end{questao}

\begin{questao}[2]{}
A fonte de corrente dependente está na periferia da malha 2 e impõe
$I_2=0{,}5I_1$. Determine $I_1$, $I_2$ e o módulo da tensão sobre essa fonte.

\circuitoq{%
\begin{circuitikz}[american,scale=.92,transform shape,line width=.8pt]
\draw (0,0) to[\fonteV,l^={$\SI{30}{\volt}$}] (0,3.2)
      to[R,l={$\SI{5}{\ohm}$}] (3.2,3.2) coordinate(t);
\draw (t) to[R,l={$\SI{5}{\ohm}$}] (3.2,0) coordinate(b) -- (0,0);
\draw (t) to[R,l={$\SI{10}{\ohm}$}] (6.4,3.2)
      to[cisource,l={$0{,}5I_1$},invert] (6.4,0) -- (b);
\draw[->,loopcol,line width=1.2pt] (1.75,1.00) arc[start angle=0,end angle=-310,radius=.46cm];
\node[loopcol,fill=white,inner sep=1pt] at (1.75,1.00){\footnotesize$I_1$};
\draw[->,loopcol,line width=1.2pt] (4.75,1.00) arc[start angle=0,end angle=-310,radius=.46cm];
\node[loopcol,fill=white,inner sep=1pt] at (4.75,1.00){\footnotesize$I_2$};
\end{circuitikz}}

\resposta{$I_1=\SI{4}{\ampere}$; $I_2=\SI{2}{\ampere}$; $|v_s|=\SI{10}{\volt}$}
\resolucao{Por estar na periferia, a fonte fixa diretamente
$I_2=0{,}5I_1$. Na malha 1,
\[
10I_1-5I_2=30.
\]
Logo $10I_1-2{,}5I_1=30$, de onde $I_1=\SI{4}{\ampere}$ e
$I_2=\SI{2}{\ampere}$. Na malha 2, as quedas resistivas totalizam
$10(2)+5(2-4)=\SI{10}{\volt}$ em módulo; a fonte de corrente sustenta essa tensão.}
\end{questao}
